\chapter{이어진 선을 찾는 법}\label{ch:find}

앞 장에서 봄베의 일을 둘로 나누었다. 위치마다 연결을 만드는 일과, 전압을 건 선에서 이어진 선을 모두 찾는
일이다. 이 장은 둘째 일을 본다. 먼저 찾은 결과가 무엇을 뜻하는지 보고, 찾는 방법 세 가지를 비교한다. 세 방법은
이 문서의 세 기계와 기준 모델에 하나씩 대응하고, 끝에서 그 대응으로 문서 전체의 지도를 그린다.

\section{전압을 걸면 무엇이 켜지는가}\label{sec:f-light}

그림~\ref{fig:c-graph}의 작은 봄베에서 케이블 \texttt{A}의 \texttt{a}번 선에 전압을 건다. ``\texttt{A}의 짝은
\texttt{A}다(짝 없음)''라는 가설을 시험하는 것이다. 전압은 연결을 따라 퍼진다. $\Scr_1$을 지나 \texttt{B}의
\texttt{b}번 선이, $\Scr_3$을 거꾸로 지나 \texttt{C}의 \texttt{c}번 선이 켜지고, 그 선들에서 다시 퍼진다. 더 퍼질
곳이 없을 때 켜진 선은 그림~\ref{fig:f-groups}의 파란 점 12개다.

\begin{figure}[ht]
\centering
\begin{tikzpicture}[x=1cm,y=-5.5mm,
    dot/.style={circle,draw,fill=white,inner sep=1.5pt},
    b/.style={fill=cblue,draw=cblue}, g/.style={fill=cgreen,draw=cgreen}, o/.style={fill=corange,draw=corange},
    lb/.style={line width=1.1pt,cblue}, lg/.style={line width=1.1pt,cgreen}, lo/.style={line width=1.1pt,corange}]
  \foreach \col/\name/\key in {0/A/A,3/B/B,6/C/C,9/A/Z} {
    \node[font=\small\bfseries] at (\col,-1.1) {\texttt{\name}};
    \foreach \l [count=\i from 0] in {a,b,c,d,e,f} { \coordinate (\key\l) at (\col,\i); }
  }
  \foreach \l [count=\i from 0] in {a,b,c,d,e,f} {
    \node[font=\ttfamily\scriptsize,anchor=east] at (-0.15,\i) {\l};
  }
  % S1: a-b b-a c-f d-e e-d f-c
  \draw[lb] (Aa) -- (Bb); \draw[lb] (Ab) -- (Ba); \draw[lb] (Ac) -- (Bf); \draw[lg] (Ad) -- (Be);
  \draw[lb] (Ae) -- (Bd); \draw[lo] (Af) -- (Bc);
  % S2: a-b b-a c-d d-c e-f f-e
  \draw[lb] (Ba) -- (Cb); \draw[lb] (Bb) -- (Ca); \draw[lo] (Bc) -- (Cd); \draw[lb] (Bd) -- (Cc);
  \draw[lg] (Be) -- (Cf); \draw[lb] (Bf) -- (Ce);
  % S3: a-c b-e c-a d-f e-b f-d
  \draw[lb] (Ca) -- (Zc); \draw[lb] (Cb) -- (Ze); \draw[lb] (Cc) -- (Za); \draw[lo] (Cd) -- (Zf);
  \draw[lb] (Ce) -- (Zb); \draw[lg] (Cf) -- (Zd);
  \foreach \k in {A,Z} { \foreach \l in {a,b,c,e} { \node[dot,b] at (\k\l) {}; }
    \node[dot,g] at (\k d) {}; \node[dot,o] at (\k f) {}; }
  \foreach \l in {a,b,d,f} { \node[dot,b] at (B\l) {}; } \node[dot,g] at (Be) {}; \node[dot,o] at (Bc) {};
  \foreach \l in {a,b,c,e} { \node[dot,b] at (C\l) {}; } \node[dot,g] at (Cf) {}; \node[dot,o] at (Cd) {};
  \draw[<-,line width=1pt] ([xshift=-1mm]Aa) ++(-0.35,0) -- ++(-0.9,0) node[left,font=\footnotesize] {전압};
  \node[note,anchor=west,align=left,text=cblue] at (9.5,0.5) {덩어리 1 (12개)\\전압을 건 선의 덩어리};
  \node[note,anchor=west,align=left,text=cgreen] at (9.5,3.0) {덩어리 2 (3개)\\\texttt{A}$\to$d, \texttt{B}$\to$e, \texttt{C}$\to$f};
  \node[note,anchor=west,align=left,text=corange] at (9.5,5.0) {덩어리 3 (3개)\\\texttt{A}$\to$f, \texttt{B}$\to$c, \texttt{C}$\to$d};
\end{tikzpicture}
\caption{그림~\ref{fig:c-graph}의 연결은 덩어리 셋으로 나뉜다. \texttt{A}의 \texttt{a}번 선에 건 전압은 파란 덩어리에만
퍼진다. 케이블 \texttt{A}에서는 \texttt{a}, \texttt{b}, \texttt{c}, \texttt{e}가 켜지고 \texttt{d}, \texttt{f}가 꺼진 채로 남는다.}
\label{fig:f-groups}
\end{figure}

이 연결 18개는 서로 만나지 않는 덩어리 셋으로 나뉜다. 전압은 건 선의 덩어리 안에서만 퍼지므로 다른 덩어리의
선은 꺼진 채로 남는다. 어느 선에 전압을 걸든 결과는 ``그 선이 속한 덩어리 전체가 켜진다''이다.

\section{켜진 선을 읽는 법: 테스트 레지스터}\label{sec:f-register}

켜진 선이 무엇을 말해 주는가? 핵심은 앞 장의 두 사실이다. 연결 하나는 추론 하나이고, 케이블 하나의 26가닥
가운데 참인 가설은 꼭 하나다.

\Why 전압을 건 가설이 참이라고 해 보자. 그러면 연결을 따라 켜진 선은 모두 그 가설에서 추론한 것이므로 모두
참이어야 한다. 그런데 그림~\ref{fig:f-groups}에서 케이블 \texttt{A}에 켜진 선이 넷이다. 한 글자의 짝이 넷일 수는
없다. 그러므로 처음 가설은 거짓이다. 연결에는 방향이 없으므로 같은 논리가 켜진 넷 모두에 적용된다. 넷 가운데
어느 것에서 출발해도 같은 덩어리가 켜지기 때문이다. 넷은 함께 서고 함께 넘어지며, 여기서는 함께 넘어졌다.
남은 것은 꺼진 \texttt{d}와 \texttt{f}다.

이제 두 경우를 나눈다.

\begin{itemize}
\item 케이블 \texttt{A}의 선이 \emph{전부} 켜지면, 짝에 대한 가설이 하나도 남지 않는다. 이 위치에서는 메뉴를 만족하는
플러그보드가 있을 수 없다. 다음 위치로 간다.
\item 하나라도 꺼진 선이 남으면, 그 가설은 아직 모순이 드러나지 않았다. 기계가 \emph{정지}한다.
\end{itemize}

그래서 봄베는 케이블 하나만 지켜본다. 이 케이블을 \term{테스트 레지스터}{test register}라 하고, 그 글자를
\term{테스트 글자}{test letter}라 한다. 판정은 한 문장이다. \emph{테스트 레지스터의 선이 전부 켜졌는가.}

표~\ref{tab:f-two}는 같은 메뉴를 두 위치에서 시험한 것이다. 두 위치는 $\Scr_3$만 다르다. 위치 1은 앞의 그림이고,
위치 2에서는 전압이 18개 선 전부에 퍼져 테스트 레지스터가 찬다.

\begin{table}[ht]
\centering
\caption{작은 봄베의 두 위치. 두 경우 모두 \texttt{A}의 \texttt{a}번 선에 전압을 걸었다.}
\label{tab:f-two}
\small
\begin{tabular}{l l l l}
\toprule
& $\Scr_3$ & 테스트 레지스터(케이블 \texttt{A}) & 판정 \\
\midrule
위치 1 & a$\leftrightarrow$c, b$\leftrightarrow$e, d$\leftrightarrow$f & a b c e 켜짐, d f 꺼짐 & 정지 \\
위치 2 & a$\leftrightarrow$c, b$\leftrightarrow$d, e$\leftrightarrow$f & 여섯 모두 켜짐 & 다음 위치로 \\
\bottomrule
\end{tabular}
\end{table}

\General 스도쿠를 풀다 막히면 ``이 칸이 3이라고 해 보자''로 시작해 귀결을 따라간다. 같은 줄에 3이 둘 나오면
가정이 틀린 것이다. 봄베는 이것을 가설 26개에 대해, 위치마다 한다. 다만 가정을 하나씩 해 보지 않는다. 가설 하나에
전압을 걸면 같은 덩어리의 가설이 한꺼번에 시험되고, 위의 판정은 그 덩어리 전체에 대한 것이다.

\Ours 실물 두 기계는 테스트 레지스터의 26선마다 \emph{감지 릴레이}를 하나씩 달아, 선에 전류가 오면 릴레이가 붙고
스스로를 붙들게 한다(\ref{sec:t-register}절). FPGA는 26비트를 ``모두 1''과 비교하고(\code{test\_q != FULL}), Rust는
정수 하나를 비교한다(\code{reg != FULL}, \ref{ch:rust-bombe}장).

\paragraph{정지는 답이 아니라 후보다.} 위치 1에서 꺼진 채 남은 것은 \texttt{d}와 \texttt{f} 둘이다. 이 작은 봄베의
메뉴는 실제로 ``\texttt{A}의 짝은 \texttt{D}, \texttt{B}는 \texttt{E}, \texttt{C}는 \texttt{F}''인 플러그보드로 만든
것이고, 그것이 초록 덩어리다. 참인 가설의 덩어리는 거짓으로 넘어갈 수 없으므로 언제나 꺼진 채 남는다. 그런데 주황
덩어리(\texttt{A}의 짝은 \texttt{F})도 거짓인데 꺼진 채 남았다. 이 덩어리 안에서는 모순이 드러나지 않았을 뿐이다.
틀린 위치에서도 우연히 모순이 드러나지 않으면 정지가 난다. 정지한 위치에서 남은 가설을 하나씩 따져 진짜를 가리는
뒷단계를 \term{체킹}{checking}이라 한다(\ref{ch:rust-check}장). 실제 봄베의 정지도 대부분 이런 거짓 정지였다.

\paragraph{모순은 루프에서 나온다.} 메뉴가 루프 없이 \texttt{A}--\texttt{B}--\texttt{C}로 끝난다고 해 보자. 그러면
케이블 \texttt{A}의 선 하나는 \texttt{B}의 선 하나, \texttt{C}의 선 하나와만 이어지고, 다시 \texttt{A}로 돌아오는 길이
없다. 케이블 \texttt{A}에 켜지는 선은 언제나 하나뿐이어서 모순이 드러나지 않는다. 모든 위치에서 정지한다. 루프가
있어야 추론이 출발한 케이블로 돌아와 ``짝이 둘''을 만든다. 해독자가 루프가 많은 크립을 고르고(그림~\ref{fig:g-menu}의
파란 루프), 루프가 지나는 글자를 테스트 글자로 삼는 까닭이다. 대각선 보드는 루프가 아닌 길에서도 케이블 사이를
되돌아오는 길을 더해 주므로, 루프가 적은 메뉴로도 모순을 만든다.

\section{방법 하나: 전류에 맡긴다}\label{sec:f-current}

\Why 이어진 선을 찾는 가장 쉬운 방법은 찾지 않는 것이다. 선이 실제 전선이고 연결이 실제 금속의 접촉이면, 전압을
거는 순간 이어진 선이 모두 같은 전압이 된다. 순서도, 단계도, ``다 찾았는가''의 확인도 없다. 676가닥이 몇 개의
연결을 거치든 거의 한순간이다. 탐색 전체를 물리가 공짜로 해 준다.

남는 일은 \emph{시간을 맞추는 것}이다. 연결이 다 자리 잡은 뒤에 전압을 걸고, 감지 릴레이가 붙을 만큼 기다린 뒤에
판정하고, 다음 위치로 연결이 바뀌기 전에 전압과 기억을 지워야 한다. 릴레이가 붙는 데에 수 밀리초가 걸리고 그
시간이 개체마다 들쭉날쭉하다는 것이 실물 기계 설계의 어려움이다(\ref{ch:relay-timing}장).

\General 물을 부어 미로의 출구를 찾는 것과 같다. 입구에 물을 부으면 물은 이어진 통로를 모두 동시에 채운다.
갈림길마다 고민하는 것은 물이 아니라 미로 속을 걷는 사람이다.

\Ours 1940년대의 봄베, 그리고 이 프로젝트의 기어식과 릴레이식이 이 방법을 쓴다. 두 기계가 다른 것은 연결을
만드는 방법(드럼을 돌리는가, 접점을 바꾸는가)과 시간을 맞추는 방법(슬롯 원판인가, 순서기인가)이고, 이어진 선을
찾는 방법은 같다.

\section{방법 둘: 한 가닥씩 따라간다}\label{sec:f-stack}

\Why 프로그램에는 전류가 없다. 선은 메모리의 비트이고, 비트는 이웃 비트가 켜졌다고 저절로 켜지지 않는다. 그래서
하나씩 따라간다. ``켜졌지만 아직 이웃을 살피지 않은 선''의 목록을 두고, 목록에서 선을 하나 꺼내 그 선과 연결된
선 가운데 꺼진 것을 켜고 목록에 넣는다. 목록이 비면 끝이다.

이 방법이 빠짐없이 찾는 이유는 간단하다. 닿을 수 있는 선은 언젠가 목록에 들어가고, 목록에 들어간 선은 언젠가
꺼내져 이웃이 살펴진다. 한 번 켜진 선은 다시 넣지 않으므로 반드시 끝난다. 그림~\ref{fig:f-groups}에서 이 방법으로
따라가면, 한 가지 순서로는 \texttt{Aa}, \texttt{Bb}, \texttt{Cc}, \texttt{Bd}, \texttt{Ae}, \texttt{Cb}, \texttt{Ba},
\texttt{Ab}, \texttt{Ce}, \texttt{Bf}, \texttt{Ac}, \texttt{Ca}의 차례로 켜진다(\texttt{Aa}는 케이블 \texttt{A}의
\texttt{a}번 선). 열두 번 꺼내고 끝난다.

\General 그림판의 ``채우기''가 이 방법이다. 한 점에서 시작해 같은 색으로 이어진 점을 모두 칠한다. 칠한 점을
목록에 넣고, 하나씩 꺼내 상하좌우를 살핀다.

\Ours Rust 기준 모델의 \code{propagate} 함수가 이 방법이다(\ref{sec:propagate}절). 컴퓨터는 한 가닥씩 따라가지만
아주 빨라서, PC 한 대가 초당 수백만 위치를 시험한다. 덧붙여, 기판을 검증하는 프로그램은 같은 물음을 또 다른 방법으로
푼다. 연결 하나를 볼 때마다 두 덩어리를 하나로 합치는 \term{합집합--찾기}{union--find}다(\ref{sec:v-uf}절).

\section{방법 셋: 모두 함께 한 걸음씩}\label{sec:f-rounds}

\Why FPGA는 게이트와 전선으로 된 회로이니 전류에 맡기면 될 것 같지만 그렇지 않다. 게이트에는 방향이 있다. 입력이
출력을 정하지 그 반대는 아니다. 방향 없는 연결을 게이트로 쓰면 ``\texttt{B}의 선은 \texttt{A}의 선으로 켜지고''와
``\texttt{A}의 선은 \texttt{B}의 선으로 켜진다''가 함께 있어야 하고, 게이트의 출력이 돌고 돌아 자기 입력이 되는
고리가 생긴다. 이 고리는 한번 켜지면 꺼지지 않아서, 다음 위치를 시험하려고 지울 방법이 없다(\ref{ch:rtl-core}장).

그래서 FPGA는 676가닥의 상태를 레지스터(기억 소자)에 담아 두고, 클럭이 한 번 뛸 때마다 \emph{켜진 선 전부가 동시에
한 걸음씩} 퍼지게 한다. 걸음마다 새로 켜진 선이 없으면 끝이다. 표~\ref{tab:f-rounds}는 그림~\ref{fig:f-groups}를
이렇게 푼 것이다.

\begin{table}[ht]
\centering
\caption{위치 1을 걸음 단위로 퍼뜨린 결과. 걸음마다 켜진 선 \emph{모두}의 이웃이 함께 켜진다.}
\label{tab:f-rounds}
\small
\begin{tabular}{c l c}
\toprule
걸음 & 새로 켜진 선 & 케이블 \texttt{A}에 켜진 선 \\
\midrule
0 & \texttt{Aa} (전압) & a \\
1 & \texttt{Bb}, \texttt{Cc} & a \\
2 & \texttt{Bd}, \texttt{Ca} & a \\
3 & \texttt{Ac}, \texttt{Ae} & a c e \\
4 & \texttt{Bf}, \texttt{Cb} & a c e \\
5 & \texttt{Ba}, \texttt{Ce} & a c e \\
6 & \texttt{Ab} & a b c e \\
7 & 없음 --- 끝 & a b c e (정지) \\
\bottomrule
\end{tabular}
\end{table}

더 바뀌지 않는 상태를 \term{고정점}{fixed point}이라 한다. 위치 1은 여섯 걸음 만에 고정점에 닿았고, 일곱째 걸음에서
바뀐 것이 없음을 확인했다. 위치 2는 아홉 걸음이 걸린다. 테스트 레지스터가 다 차면 더 볼 것이 없으므로 거기서
끝낸다.

\General 산불이 번지는 것을 한 시간마다 지도에 칠한다고 생각하면 된다. 매시간, 지난 시간까지 탄 곳의 이웃이 모두
함께 탄다. 새로 칠할 곳이 없으면 불이 다 번진 것이다. 몇 시간이 걸리는지는 지형이 정한다.

\Ours 이 방법의 대가는 두 가지다. 첫째, 걸음 수만큼 클럭이 든다. 이 프로젝트의 FPGA 엔진은 위치 하나에 약 12--17클럭을
쓴다. 둘째, 한 걸음에 \emph{모든} 연결을 동시에 처리하므로, 연결 수백 개마다 회로가 따로 있어야 한다. 이 회로를 어떻게
적느냐에 따라 크기가 약 열 배 달라졌다. 첫 구현은 칩 용량의 약 세 배였다(\ref{ch:lesson}장). FPGA 쪽 설계의 많은
부분이 ``전류가 공짜로 하던 일을 회로로 얼마나 싸게 짓는가''이다.

\section{셋이 같은 답을 내는 까닭}\label{sec:f-same}

세 방법은 따라가는 순서가 전혀 다르다. 전류에는 순서가 없고, 목록은 한 가닥씩, 걸음은 한 겹씩 간다. 그래도 결과는
같다. 구하는 것이 ``전압을 건 선의 덩어리''이고 덩어리는 하나로 정해져 있기 때문이다. 어떤 순서로 따라가든 닿을 수
있는 선에는 언젠가 닿고, 닿을 수 없는 선에는 끝내 닿지 않는다. 켜진 선이 다시 꺼지는 일도 없다.

이 사실이 이 프로젝트의 검증을 받친다. 세 기계와 Rust 모델은 같은 메뉴에서 같은 정지를 내야 하고, 그래서 정지
목록을 하나하나 맞대어 볼 수 있다. Rust 모델을 정답표로 삼아 FPGA 회로를 맞대는 방법은 \ref{ch:verify}장에서, 릴레이
기판을 맞대는 방법은 \ref{ch:boardverify}장에서 본다.

\section{무엇이 공짜이고 무엇이 비싼가}\label{sec:f-cost}

두 가지 일을 다시 놓고 보면, 실물 기계와 FPGA는 비싼 쪽이 정반대다.

\begin{itemize}
\item \textbf{실물 기계}에서 이어진 선을 찾는 일은 공짜다. 비싼 것은 연결을 바꾸는 일이다. 드럼이 한 칸 돌거나
릴레이가 넘어가는 데 밀리초가 들고, 그 사이 판정 시각을 맞춰야 한다.
\item \textbf{FPGA}에서 연결을 바꾸는 일은 싸다. 표에서 다른 칸을 읽으면 되고 나노초면 된다. 비싼 것은 이어진 선을
찾는 일이다. 걸음을 되풀이해야 하고, 그 걸음을 위한 회로가 칩을 차지한다.
\end{itemize}

표~\ref{tab:f-speed}는 그 결과다. 실물 두 기계는 1940년대 봄베와 같은 규모의 속도이고, FPGA는 약 10만 배 빠르다.

\begin{table}[ht]
\centering
\caption{위치 하나에 드는 시간. 이 프로젝트의 값은 모두 설계값이며 실측 전이다(보고서 6장의 표에서 옮김).}
\label{tab:f-speed}
\small
\begin{tabular}{l r l}
\toprule
기계 & 위치당 & 시간을 정하는 것 \\
\midrule
영국 봄베(1940년대) & 35 ms & 드럼의 회전 \\
기어식 & 44 ms (추정) & 드럼의 회전 \\
릴레이식 & 63 ms (계산) & 릴레이가 넘어가고 붙는 시간 \\
FPGA (50 MHz) & 0.3 $\mu$s (추정) & 퍼뜨리는 걸음 수 \\
\bottomrule
\end{tabular}
\end{table}

\section{이 문서의 지도}\label{sec:f-map}

표~\ref{tab:f-map}가 이 문서의 지도다. 가로줄이 기계 하나이고, 세로줄이 그 기계가 각 일을 하는 수단과 그것을 다루는
장이다.

\begin{table}[ht]
\centering
\caption{세 기계와 기준 모델이 봄베의 일을 하는 수단.}
\label{tab:f-map}
\small
\begin{tabularx}{\linewidth}{l >{\raggedright\arraybackslash}X >{\raggedright\arraybackslash}X >{\raggedright\arraybackslash}X >{\raggedright\arraybackslash}X}
\toprule
& 선 676가닥 & 연결을 만든다 & 이어진 선을 찾는다 & 결과를 읽고 박자를 맞춘다 \\
\midrule
Rust 모델 & 정수 26개의 비트 (\ref{ch:rust-bombe}장) & 미리 계산한 치환표 & 목록으로 한 가닥씩 & \code{reg != FULL}, 차례대로 \\
기어식 & 전선 & 드럼이 돈다 (\ref{ch:50-gear-mechanics}장) & 전류 & 감지 릴레이 (\ref{ch:relay-timing}장), 슬롯 원판 (\ref{sec:t-gear}절) \\
릴레이식 & 전선 & 접점이 밀기를 고른다 (\ref{ch:relay-machine}장) & 전류 & 감지 릴레이, 순서기 (\ref{ch:relay-sequencer}장) \\
FPGA & 레지스터의 비트 & 메모리의 치환표 (\ref{ch:rtl-scrambler}장, \ref{ch:fpga-engine}장) & 클럭마다 한 걸음 (\ref{ch:rtl-core}장, \ref{ch:lesson}장) & \code{test\_q != FULL}, 클럭 (\ref{ch:sync}장) \\
\bottomrule
\end{tabularx}
\end{table}

이 문서는 이 표를 위에서 아래로 읽는다.

\begin{itemize}
\item \textbf{B부}는 Rust 모델이다. 모든 기계가 내야 할 답을 정하는 정답표이고, 위의 일들을 가장 읽기 쉬운 형태로 보여 준다.
\item \textbf{C부}는 실물 두 기계가 함께 쓰는 부품이다. 전기의 기초, 릴레이, 그리고 결과를 읽고 박자를 맞추는 시간 규칙과 테스트 레지스터.
\item \textbf{D부}는 기어식, \textbf{E부}는 릴레이식이다. 연결을 만드는 방법과 박자를 맞추는 방법(슬롯 원판, 순서기)이 어떻게 다른지가 중심이다.
\item \textbf{F부}는 릴레이식 회로를 FPGA로 옮기면서 무엇이 바뀌는지를 개념별로 보고, \textbf{G부}는 그 FPGA의 코드를 읽는다.
\item \textbf{H부}는 실물을 만들고 검증하는 도구와 전원·안전, \textbf{I부}는 기계를 부리는 소프트웨어다.
\end{itemize}

\section{정리}

\begin{itemize}
\item 연결은 선들을 덩어리로 나눈다. 전압을 건 선의 덩어리 전체가 그 가설의 귀결이다.
\item 테스트 레지스터의 선이 전부 켜지면 가설이 하나도 남지 않으므로 다음 위치로 간다. 하나라도 꺼져 있으면 정지한다.
정지는 후보이고, 체킹이 가린다. 모순은 루프(와 대각선 보드)에서 나온다.
\item 이어진 선을 찾는 방법은 셋이다. 전류에 맡기기(실물 기계), 목록으로 한 가닥씩(Rust), 모두 함께 한 걸음씩(FPGA).
덩어리는 하나로 정해져 있으므로 셋은 같은 답을 낸다.
\item 실물 기계는 찾기가 공짜이고 연결을 바꾸기가 비싸다. FPGA는 그 반대다.
\end{itemize}
